\section{Rational lift tests on the original tied family}

% BEGIN T76R1
\paragraph{T76R1. The core, the two marked primes and injective evaluation.}
Fix the original three-stage datum $R_3$, face $F_3$, geometric
core $T$, weights $\kappa$, vectors $v_0$ and base records $O_i$.
Assume $\mathrm{TiesB}(R_3,F_3)$ and
$\mathrm{TiesG}(T,R_3,\kappa,v_0,O)$. Put $R=R_3.\mathrm{toRealDat}$,
$F=F_3.\mathrm{lv2}$ and take the weight and frame tie $w$
in the latter hypothesis. The weak characterization gives the level,
group, family and model laws at $F$. At each original index,
apply the conditional chain identities for the two intertwiners,
the oscillator maps and theta to their existence witnesses in
$\mathrm{TiesB.ex}$. Do the same for the point oscillator,
$\mathrm{thF}$ and $\mathrm{thG}$, retaining both the action on vectors
and the second clause for $\mathrm{thF}$. This supplies the chain
and point-theta fields and hence all six fields of $\mathrm{Hk3.Core}(R,F)$.
The tuple laws in the same weak characterization give the off-place
and used-place character identities.

The base records determine the two marked primes of $T$ in their
specified order. Since $2\le D.R.m$, take the row $i=0$ and its
unique base tuple. For $k=0,1$, choose by lying over a prime $P$ of
$\mathcal O_E$ above the nonzero marked prime $R.C.P.\mathrm{pb}(k)$
of $\mathcal O_K$. The level law puts $P$ in $R.P_g$. The tuple half
pin then puts $T.Q.l.M_b$ in $P$. By the geometric half law of
$w.\mathrm{laws}$ this ideal is the product of the extensions of
the marked primes of $T$; hence $P$ contains one such factor and
contracts to $T.\mathrm{pb}(k')$. The contraction contains both
rational primes $g_k$ and $g_{k'}$. If $k\ne k'$, these are
coprime, an impossibility in a proper ideal. Thus
$T.\mathrm{pb}=R.C.P.\mathrm{pb}$.

To prove $\mathrm{EvInj}(T,R,v_0)$, let $x,y$ give the same
evaluation at one of the prescribed $v=v_0(i,\sigma)$. The
multiplicity map $U=e(x-y)$ vanishes at $v$. Each local component of
$v$ has norm one by the chosen extremal-vector condition of $w$.
The product $\prod_\tau\langle v_\tau,v_\tau\rangle=1$, so the
defining multiplicity-inner-product identity at $(v,v)$ gives
$\mathrm{ipM}(U,U)=0$ from the zero integrand. In the alternative
defining case with no inner-product witness, $\mathrm{ipM}$ is
zero by definition. The target model law identifies
$\langle x-y,x-y\rangle$ with the prescribed multiple of this
zero, so $x=y$. This works at every required embedding.
% END T76R1

% BEGIN T76R2
\paragraph{T76R2. All tuple records over the field of the pinned row.}
Let $L=D.\mathrm{what}_2.\mathrm{LL}$, $\sigma=D.\mathrm{what}_2.\mathrm{embLL}$
and $M=\sigma(L)\subset\mathbb C$. The base pin at each $i$ gives
$O_i.c'.F=M$. Fix this coefficient record together with its map
$O_i.j_L$ and its inclusion of the lower coefficient field.
The level law gives finite order, triviality on the base field,
tameness and the ramification and residue identities for the tuple
characters. The family law gives surjectivity on the allowed primes
and the specified prime values of the twisted family through its
coefficient tie. Consequently RTG2 gives $\mathrm{XiVal}(R)$.
In particular, off the enlarged conjugation-stable exceptional set
one takes the quotient of the tuple and base values at a moved
prime, and at a fixed prime the square of the value is one.
RTG2 then obtains values at the exceptional places by principal
ideles and passes to the adelic character by continuity and its
finite image.

The two clauses of $\mathrm{TiesB.guard}$ give agreement of the
selected $K$-prime for a common rational residue prime, and exclusion
of the original bad integer at all selected $E$-primes. Use
RTG10 and the ordered prime equality of T76R1 to form the prescribed
half ideal and cover level for each tuple. The roots in
$R.\mathrm{fam\_roots}$ come from the same family law. Together
with the primitive lower-level root they supply the upper-level
roots inside $M$. Apply D74U5 with $w.\mathrm{laws}$, the level
law, these roots, $\mathrm{XiVal}$, and $\mathrm{TiesG.cover}$.
The last field supplies the arithmetic function-field condition
on the actual constructed higher cover. This gives
$\mathrm{CoversAt}$ on $O_i.c'$ with its original inclusion.
RTG12 now gives $\mathrm{PinsOverL}$, including the placed
coefficient vectors, symbol factors and residue functions for
all tuples. Its base coefficient record and level map agree
with those of $O_i$ as in (RTG.17).
% END T76R2

% BEGIN T76R3
\paragraph{T76R3. The phase-correct family of rational vectors.}
For each $i$, take the coefficient witness
$(\mathrm{gramR},c_{0R},c_{0L})$ from the lift pin of $O_i$ and
retain $O_i.\mathrm{lev}$. Retie the record of T76R2 to that
witness by RTG4--RTG5. Its pins and coefficient field survive.
The core and the two tuple laws of T76R1, $w$, the equality of
the marked primes, and $\mathrm{TiesG.ioExt}$ are exactly the outer
inputs to P74I7. With the fixed coefficient witness and level map
it supplies a single nowhere-zero $\lambda_i$ and a single level
element before either record or its character constant is tested.
T73C5 gives class-sum existence on every such pinned record,
including the one after any allowed $\mathrm{scaleXi}$. Its field
inclusion is equality here and its prime agreement is T76R1.

The rank inequality $8D.R.f\,D.R.r+4\le D.R.nV$ in
$\mathrm{TiesB.nV}$, the core, $\mathrm{ioExt}$, and the weight,
vector, frame and algebra laws of $w$ give F76B2 on each $O_i$.
Thus every tuple class sum is a finite linear combination of the
base class sums over their common field. Reading both sides through
the same $\lambda_i$ and level element gives RTG16's
$\mathrm{LiftsSpan}$. Use the first clause of
$\mathrm{TiesG.oneMap}$ for the base map. Evaluating
$\mathrm{TiesG.phase}$ at $\mathrm{TiesG.unif}$ fixes the phase
\[
       p_i=\operatorname{cj}_{s}(R_3.\mathrm{phs}_i),\qquad
       s=\mathrm{twAt}(\sigma,A_3.A.v_0).
\]
When a base lift is nonzero, RTG14 and RTG17 put the ratio of the
two character constants in $M$. Retie by this one scalar and reuse
the class-sum and point identities on the resulting record. If
every base lift is zero, the span already gives zero for every
tuple lift. RTG18 gives $\mathrm{OneMapOverL}$ with $p_i$ in
either case. Finally $w$ and $\mathrm{TiesG.unif}$ give its
evaluation law by RTG13, and T76R1 supplies injectivity. RTG19--RTG20
therefore give, for every $i,t,x$, a rational coordinate vector
$y\in L^{R.nW_i}$ with
\[
 R.FW_i.j(\mathrm{sg0})y
       =p_i\,R.\mathrm{thG}_{i,t}(R.jV_{i,t}x).
\]
This is $\mathrm{RatG}$ for the prescribed phase.
% END T76R3

% BEGIN T76R4
\paragraph{T76R4. The tests and lattices on the same three-stage object.}
For $X=R_3.\mathrm{toCx3}(F_3)$, the original
$\mathrm{TiesB}$ and $\mathrm{TiesG}$ hypotheses imply
\[
  \mathrm{LiftTest}(X),\qquad
  \mathrm{LiftTestB}(X,b),\qquad X.\mathrm{MLat}(b)
    \quad (b=\mathrm{false},\mathrm{true}).
\]
Indeed RTG21 twists the common independent spanning coordinates
from $\mathrm{TiesG.unif}$ to obtain separating linear tests.
T76R3 makes their values on lifted rational inputs rational in
the prescribed embedding. The rational basis of the good subspace
at each level gives the last test field there. The tie $w$ also
gives $\mathrm{LawsJ}$ and $\mathrm{KapIs}$, so M76L2 and the
coordinate transport of RTG22--RTG24 give the lattice on both
levels of $X$. All maps and bases here descend from the same
$R_3,F_3$. In particular the specialization $b=\mathrm{false}$
is the native conclusion of
$\texttt{liftTests\_and\_mLat\_of\_ties}$ on its own hypotheses.
Producing $T,R_3,F_3,\kappa,v_0,O$ with the simultaneous ties
required by the outer existence statements is a further obligation.
% END T76R4
